% !TeX root = ../main.tex
\begin{frame}{01 研究背景——血管支架的功能}
  \begin{SlideCanvas}
    \BodyAt{55}{84}{850}{支架以恢复血流、维持管腔通畅为共同目标，设计要求可归纳为\BICEmph{几何覆盖}、\BICEmph{径向支撑}与\BICEmph{界面适配}三个维度。}
    \PicAt{55}{161}{535}{295}{assets/gupta2026_fig1.jpg}
    \GroupHeadingAt{610}{153}{290}{支架功能需求}
    \draw[BICNavy,line width=.8pt] (610,190)--(900,190);
    \BulletPointAt{630}{211}{270}{几何覆盖}
    \DetailAt{630}{244}{270}{覆盖目标病变段，控制展开长度与端部位置。}
    \BulletPointAt{630}{300}{270}{径向支撑}
    \DetailAt{630}{333}{270}{恢复并维持所需管腔，抵抗外部压缩与回弹。}
    \BulletPointAt{630}{389}{270}{界面适配}
    \DetailAt{630}{422}{270}{贴合血管壁，协调支架与组织的局部受力。}
    \Sources{图源：\RefGupta，图1。功能与临床机制：\RefSCAI。}
  \end{SlideCanvas}
  \BICConclusionBand{几何覆盖、径向支撑与界面适配共同服务于\BICBandEmph{管腔通畅}}
  \note{三项要求是共同目标下的并列设计维度，不强行排列为依次发生的过程。几何覆盖处理病变段和端部位置，径向支撑处理管腔尺寸和机械维持，界面适配处理贴壁和局部组织受力。图1(a)为狭窄与植入示意，(b)(c)为不同结构变形。不同血管部位和部署方式的量化要求需按对应工况解释。}
\end{frame}
